% ============================================================
% RESUMEN / ABSTRACT
% ============================================================

\chapter*{Resumen}
\addcontentsline{toc}{chapter}{Resumen}

\section*{Título del trabajo}
Predicción de fibrilación auricular paroxística a partir del electrocardiograma en ritmo
sinusal mediante el análisis con redes neuronales profundas

\section*{English Title}
{\itshape Prediction of Paroxysmal Atrial Fibrillation from the Sinus Rhythm
Electrocardiogram Using Deep Neural Networks}

\section*{RESUMEN}
\selectlanguage{spanish}

La Fibrilación Auricular Paroxística (PAF) es una de las arritmias más comunes, caracterizada por latidos irregulares en la aurícula.
Debido a su naturaleza intermitente que la caracteriza, su diagnóstico temprano es complejo, es por ello que la predicción de dichos 
episodios a paritr de señales ECG en ritmo sinusal es de especial interés. En este trabajo, se ha abordado este desafío.

Una de las limitaciones que suelen tener trabajos de este tipo es la falta de generalización de los modelos, ya que si son entrenados
en una base de datos específica, de unas características concretas, no suelen generalizar bien a otras bases de datos. 
En este trabajo, se han consolidado cinco bases de datos públicas de PhysioNet para abordar esta limitación. 

Se ha diseñado un preprocesamiento que extrae ventanas inmediatamente anteriores al inicio de un episodio de fibrilación auricular, 
y se extraen nueve variables de variabilidad del ritmo cardíaco (HRV). Se han implementado cuatro arquitecturas distintas de redes neuronales:
ResNet 1D (baseline), SEResNet 1D (con atención de canales), un híbrido CNN-Transformer y Mamba 1D
(\textit{selective state space model}). Para evitar la filtración de datos, se hizo validación cruzada, agrupando estrictamente por sujeto
(5-Fold GroupKFold).

Este trabajo estudia la viabilidad de un sistema generalizable inter-sujeto
para la monitorización de alerta temprana de episodios de Fibrilación Auricular Paroxística.

\vspace{0.8cm}
\noindent\textbf{Palabras clave:} Fibrilación Auricular Paroxística, Predicción Inminente, Procesamiento de ECG, Aprendizaje Profundo, Variabilidad del Ritmo Cardíaco (HRV), Mamba.

\vspace{1cm}

\section*{ABSTRACT}
\selectlanguage{english}

Paroxysmal Atrial Fibrillation (PAF) is one of the most common arrhythmias, characterized by irregular heartbeats in the atrium. Due to its intermittent nature that characterizes it, its early diagnosis is complex; that is why the prediction of such episodes from ECG signals in sinus rhythm is of special interest. In this work, this challenge has been addressed.

One of the limitations that works of this type usually have is the lack of generalization of the models, since if they are trained on a specific database, with concrete characteristics, they do not usually generalize well to other databases. In this work, five public databases from PhysioNet have been consolidated to address this limitation.

A preprocessing has been designed that extracts windows immediately prior to the onset of an atrial fibrillation episode, and nine heart rate variability (HRV) variables are extracted. Four different neural network architectures have been implemented: 1D ResNet (baseline), 1D SEResNet (with channel attention), a hybrid CNN-Transformer, and 1D Mamba (\textit{selective state space model}). To avoid data leakage, cross-validation was performed, grouping strictly by subject (5-Fold GroupKFold).

This work studies the feasibility of an inter-subject generalizable system for early-warning monitoring of Paroxysmal Atrial Fibrillation episodes.

\vspace{0.8cm}
\noindent\textbf{Keywords:} Paroxysmal Atrial Fibrillation, Imminent Prediction, ECG Signal Processing, Deep Learning, Heart Rate Variability (HRV), Mamba.

\selectlanguage{spanish} % Restaurar el idioma principal

